\chapter{Verifikation adaptiver stationärer Wärme}\label{vv:adaptive}
Die glatte Referenz auf dem Einheitswürfel ist
\(T=300+\sum_{j=1}^3 x_j(1-x_j)\), \(Q=6\) bei \(k=1\).
Die lokalisierte Referenz lautet \(T=300+\exp(-80(x-0.5)^2)\),
\(Q=(160-25600(x-0.5)^2)\exp(-80(x-0.5)^2)\).
Alle Randtemperaturen sind exakt vorgegeben. Der tatsächliche Energiefehler
\(\sqrt{\int_\Omega |\nabla(T-T_h)|^2\,dV}\) wird unabhängig vom Residualindikator
mit Gauss-Duffy-Regeln Ordnung 6 und 7 integriert; relative Parität 2e-4 ist Gate.
Tet4-Lastquadratur bleibt Grad 2 und trägt im Gaußfall ebenfalls zum Fehler bei.

Je sieben adaptive und uniforme Stufen starten am identischen Netz mit 125 DOFs.
Der glatte adaptive Fehler sinkt von 0.25 auf 0.07753 sqrt(W K), uniform auf 0.04724.
Lokalisiert sinkt er adaptiv von 1.96263 auf 1.34228. Die nach Pilotexploration
und vor dem finalen Wiederholungslauf fixierte Grenze 1.5 sqrt(W K) wird adaptiv
mit 514, uniform mit 729 DOFs erstmals erreicht. 29.5 Prozent weniger DOFs sind
fallspezifischer Nutzen; eine blind vor Pilotversuchen gewählte Grenze oder
allgemeine Laufzeitgarantie wird nicht behauptet. Effektivität bleibt heuristisch.
Uniform verfeinert dieselbe Kantenstrategie mit Markierungsanteil 1.

Regressionen prüfen vier oder mehr Stufen, Fehlerabnahme und festen DOF-Vergleich.
Weitere analytische Fälle prüfen eine Wärmeleitfähigkeitsgrenze mit stetigem
Fluss und gemischte Fluss-/Robin-/Dirichlet-Ränder. Topologie, Gruppen, Bilanzen,
identische Records und alle Abbruchgründe sind eigene Gates. Kein experimenteller
Adaptivitätsnachweis und keine transiente oder gekoppelte Freigabe.
